\subsection{函数逆元的导数}

命题 4. 设 E 是 R 上具有单位元的完备赋范代数，设 f 是定义在区间 I ⊂ R 上、取值于 E 中、且在点 \(x_{0} \in I\) 可导的函数。若 \(\mathbf{y}_{0} = \mathbf{f}(x_{0})\) 在 E 中可逆 \(^{2}\) ，则函数 \(x \mapsto (\mathbf{f}(x))^{-1}\) 在 \(x_{0}\) 的（相对于 I 的）一个邻域上有定义，且具有导数，等于 \(-\left(\mathbf{f}(x_{0})\right)^{-1}\mathbf{f}'(x_{0})\left(\mathbf{f}(x_{0})\right)^{-1}\) （在 \(x_{0}\) 处）。

事实上，E 中可逆元所成的集是一个开集，函数 \(y \mapsto y^{-1}\) 在其上连续（Gen.~Top., IX, p.~178）；由于 f 在 \(x_{0}\) （相对于 I）连续，\(\left(\mathbf{f}(x)\right)^{-1}\) 在 \(x_{0}\) 的一个邻域上有定义，且我们有

\[
\left(\mathbf {f} (x)\right) ^ {- 1} - \left(\mathbf {f} \left(x _ {0}\right)\right) ^ {- 1} = \left(\mathbf {f} (x)\right) ^ {- 1} \left(\mathbf {f} \left(x _ {0}\right) - \mathbf {f} (x)\right) \left(\mathbf {f} \left(x _ {0}\right)\right) ^ {- 1}.
\]

于是命题由 \(y^{-1}\) 在 \(y_{0}\) 的一个邻域上的连续性以及 xy 在 \(E \times E\) 上的连续性得出。

例。1) 最重要的特例是 E 为域 R 或 C 之一的情形：若 f 是实值或复值函数，在点 \(x_{0}\) 可导且 \(f(x_{0}) \neq 0\) ，则 1/f 具有导数，等于 \(-f'(x_{0})/(f(x_{0}))^{2}\) （在 \(x_{0}\) 处）。

\begin{enumerate}
\def\labelenumi{\arabic{enumi})}
\setcounter{enumi}{1}
\tightlist
\item
  若 \(U = (\alpha_{ij}(x))\) 是 n 阶方阵，在 \(x_{0}\) 可导且在该点可逆，则 \(U^{-1}\) 具有导数，等于 \(-U^{-1}U'U^{-1}\) （在 \(x_{0}\) 处）。
\end{enumerate}

